\clearpage
\section{Frozen Readers and Clean Answer Discrimination in LLaVA}
\label{app:validation-readout}

\subsection{Prospective validation allocation and qualification}

The validation screen asks which fixed clinical directions combine
controlled readability with an informative clean answer margin before
testing their write efficacy. It reuses the accepted LLaVA-1.5-7B snapshot,
native image processor, final visual block, and six train-fitted clinical
directions. The clinical questions are the six prompts in
Figure~\ref{fig:prompt-score}. The scoring interface takes the maximum
first-answer logit among accepted singleton yes variants minus the maximum
among no variants, retaining case and leading-space variants. Its sigmoid
is a model score rather than a calibrated disease probability.

Allocation uses the original manifest's validation split, excluding every
patient represented in the sixteen implementation-preflight images.
Remaining patient identifiers are sorted and permuted with seed 20260913.
The first 700 patients form the reader and capability cohort, and the next
100 are reserved for conditional write evaluation. Each patient contributes
its lexicographically smallest validation image identifier. This allocation
precedes inspection of labels and responses. The two cohorts are
patient-disjoint, as are the original training, validation, and test splits.
All six questions use the same 700 calibration images.

Clinical reader scores use the stored projection, training centering and
scaling, and logistic coefficients. Twenty recurring-type random-label
controls are fitted on the original training split and shared across the
six comparisons. No reader or control parameter is fitted on validation.
Selectivity, defined in Section~\ref{sec:protocol}, subtracts mean control
AUROC from clinical AUROC, with each control evaluated against its own type
labels. The same 2,000 patient-bootstrap draws, seeded with 20260914, serve
every clinical reader, every control, and clean answer capability.

Reader qualification requires a positive 2.5th-percentile selectivity bound.
Clean capability requires an answer-margin AUROC fifth-percentile bound
strictly above 0.5. Both require at least ten patients per label class and
1,900 valid draws. A question enters the write screen only when both
criteria pass. These prospectively fixed thresholds define an exploratory
validation policy; their percentile intervals are not selection-adjusted
confirmatory intervals across the six questions.

\subsection{Controlled reader availability}

Table~\ref{tab:validation-readers} presents all six frozen-reader comparisons.
Every question meets the class-count requirement, and every reader
comparison yields all 2,000 valid joint draws.

\begin{table}[H]
\centering
\caption{Frozen clinical-reader qualification on validation patients.}
\label{tab:validation-readers}
\resizebox{\ifdim\width>\linewidth\linewidth\else\width\fi}{!}{%
\begin{tabular}{lrrrrl}
\toprule
Question & Positive & Negative & Selectivity & 95\% interval & Qualifies \\
\midrule
Effusion & 34 & 666 & 0.1438 & $[0.0658,0.2163]$ & Yes \\
Atelectasis & 43 & 657 & 0.0472 & $[-0.0351,0.1198]$ & No \\
Pneumothorax & 12 & 688 & 0.0362 & $[-0.1271,0.1731]$ & No \\
Cardiomegaly & 25 & 675 & 0.0889 & $[0.0002,0.1723]$ & Yes \\
Mass & 39 & 661 & 0.0735 & $[-0.0023,0.1473]$ & No \\
Nodule & 48 & 652 & 0.0037 & $[-0.0833,0.0865]$ & No \\
\bottomrule
\end{tabular}%
}
\end{table}

Effusion and Cardiomegaly satisfy the registered reader screen. Effusion's
lower bound is 0.065803, whereas Cardiomegaly's is 0.000164 before rounding.
The latter sits close to the qualification threshold. The other four
intervals span zero. Effusion consequently provides the more stable
screening anchor for a subsequent independent diagnostic of answer
measurement; that choice remains a validation-stage selection.

\clearpage
\subsection{Clean capability and the conditional write stage}

Table~\ref{tab:validation-capability} reports clean answer-margin AUROC on
the same patients. Its two-sided 95\% intervals are descriptive; the
separate one-sided 95\% lower bounds implement the capability rule. All six
questions again have 2,000 valid draws. The run measures 4,200 clean
image--question outcomes in 366 seconds on one A100.

\begin{table}[H]
\centering
\caption{Clean answer discrimination across the six validation questions.}
\label{tab:validation-capability}
\resizebox{\ifdim\width>\linewidth\linewidth\else\width\fi}{!}{%
\begin{tabular}{lrrrl}
\toprule
Question & AUROC & Descriptive 95\% interval & One-sided 95\% lower bound & Qualifies \\
\midrule
Effusion & 0.4548 & $[0.3546,0.5555]$ & 0.3722 & No \\
Atelectasis & 0.5275 & $[0.4473,0.6111]$ & 0.4609 & No \\
Pneumothorax & 0.5122 & $[0.3785,0.6400]$ & 0.4028 & No \\
Cardiomegaly & 0.5232 & $[0.4034,0.6375]$ & 0.4240 & No \\
Mass & 0.5518 & $[0.4584,0.6514]$ & 0.4723 & No \\
Nodule & 0.5536 & $[0.4601,0.6482]$ & 0.4731 & No \\
\bottomrule
\end{tabular}%
}
\end{table}

The largest capability lower bound is 0.473107, so no question qualifies
jointly for reader availability and clean capability. The reserved
100-patient write cohort remains untouched. Its planned bidirectional
clinical, random, and sham comparisons therefore contribute no efficacy
estimate. This is the registered outcome of the conditional screen.

Together, the tables separate controlled direction-level information from
its reliable expression through this fixed answer margin. The reader and
capability criteria test different quantities; their qualification pattern
is not a direct statistical comparison between the measurements. Prompt
wording, answer verbalizers, and image dependence remain candidate sources
of the separation. A subsequent diagnostic can compare these factors while
independently replicating the frozen reader, before advancing to clinical
write competition.

\FloatBarrier
